% ----------------------------------------------------------
% Trabalhos Relacionados
% ----------------------------------------------------------
\chapter{Trabalhos Relacionados}

O desenvolvimento de plataformas robóticas móveis para fins educacionais e de pesquisa tem ganhado destaque nos últimos anos, em virtude da crescente popularização de computadores de placa única de baixo custo e do amadurecimento de frameworks de software abertos como o ROS.

No âmbito acadêmico, \citeonline{vilanova2022aeronave} propôs o desenvolvimento de uma plataforma robótica aérea baseada na aeronave DJI® Ryze Tello integrada ao ROS 2, voltada para práticas educacionais de robótica móvel no curso de Engenharia Mecânica do IFPI. A pesquisa evidenciou que o uso do ROS 2 em contextos educacionais favorece a compreensão dos princípios de arquitetura modular e comunicação entre nós, permitindo que o estudante vivencie o ciclo completo de desenvolvimento robótico.

Em consonância com essa proposta, \citeonline{Dantas_2022} desenvolveu uma plataforma robótica móvel adaptada a partir do aspirador Roomba® 614 iRobot, visando o ensino de robótica para o curso de Engenharia Mecânica do IFPI. O projeto foi estruturado em torno da integração com o ambiente ROS 2, viabilizando a comunicação e o controle do robô. Ademais, o trabalho introduziu uma base impressa em 3D projetada para acoplar sensores adicionais, garantindo versatilidade para diversas práticas experimentais.

O presente trabalho diferencia-se ao propor uma plataforma robótica seguidora de linha, também baseada no Roomba® 614, mas com foco na implementação de um controlador PID para navegação autônoma e na utilização do Raspberry Pi 4 Model B como controlador principal, integrando controle, sensoriamento e comunicação em uma única unidade de processamento.
